\section{Method}
\label{sec:method}

% NOTE (this revision pass): added two algorithm boxes (SELECT_TOKENS and
% CANDIDATE_DELTA/COMMIT) replacing the old prose-only description of the
% lazy-greedy loop; added a new "Precomputation" subsection with detail
% mined directly from tokenizer.py's own code comments (build_out_adjacency,
% build_coverage_transition, the multiprocessing.Pool worker-state pattern);
% converted the metrics list and the baselines list into tables
% (Table~\ref{tab:metrics} and Table~\ref{tab:baselines}); simplified
% sentence structure and added transitions between subsections. No claims
% or numbers changed from the previous version -- this is a detail/structure
% pass grounded in the actual source, not new content.

\subsection{Candidate pool and depth budget}
\label{sec:method:candidates}

A single depth budget $D$ plays three roles in this method, all controlled
by one hyperparameter, \texttt{TokenizerParam.max\_dist\_candidate} in the
implementation: it bounds how far a candidate token may be from the mapped
concepts it can help represent, it bounds how deep a context tree
(Section~\ref{sec:method:tokenizer}) may grow, and it bounds the horizon
over which the coverage score below is computed. Concretely, the candidate
pool is
\[
V_D = M \cup \{\, u \in V : u \text{ is reachable from some } c \in M
\text{ by at most } D \text{ relationships} \,\},
\]
obtained in practice as the union of the $D$-hop ego-graph (following
outgoing edges) around every mapped concept -- see Table~\ref{tab:snomed-stats}
for its current size. Token selection (Section~\ref{sec:method:selection})
only ever considers $t \in V_D$: a candidate that no mapped concept can
reach within $D$ hops has zero possible contribution to the objective below
and is never evaluated.

The current value, $D=9$, is not an arbitrary round number: it covers
98.4\% of concept depths observed in the candidate subgraph, while staying
well below the longest path length in the graph ($30$, driven by a handful
of outlier chains). Concepts only reachable beyond depth $9$ tend to be
generic SNOMED "grouper" categories (e.g.\ \emph{Wound}, \emph{Inspection
(procedure)}) rather than clinically specific concepts, so little is lost
by not extending the budget further. \doubt{Using one shared $D$ for all
three roles (candidate radius, tree depth, coverage horizon) is a
simplifying design choice worth calling out explicitly to a reviewer -- is
that coupling intentional/load-bearing, or a current implementation
convenience that a future version might relax into three independent
hyperparameters?}

\subsection{Greedy token selection}
\label{sec:method:selection}

\begin{quote}\small\itshape
\begin{itemize}
\item Greedy token selection based on semantic coverage with lambda as distance penality (add illustration)
\end{itemize}
\end{quote}

Given the candidate pool $V_D$, we now need a rule for choosing which $k$ of
its members become tokens. For $u \in V$, $T \subseteq V$, and horizon
$h \in \{0, \dots, D\}$, let $S_h(u, T) \in [0, 1]$ denote how well $u$ is
covered by $T$ when at most $h$ further relationships may be traversed. The
base case is
\[
S_0(u, T) = \mathbf{1}[u \in T],
\]
and, extending the recursion with a per-hop distance penalty $\lambda \in
(0, 1]$ \notetoself{confirmed by you: this $\lambda$ term is your own
addition on top of the original PDF draft, which only has the undiscounted
($\lambda{=}1$) version of Eq.~\eqref{eq:coverage} -- well spotted catching
that gap.},
\begin{equation}
S_h(u, T) =
\begin{cases}
1, & u \in T, \\[2pt]
0, & u \notin T \text{ and } d^+(u) = 0, \\[2pt]
\dfrac{\lambda}{d^+(u)} \displaystyle\sum_{(u,r,v) \in \mathrm{Out}(u)} S_{h-1}(v, T), & \text{otherwise.}
\end{cases}
\label{eq:coverage}
\end{equation}
In words: a token is absorbing (its coverage is $1$ at every remaining
horizon), and every other node's coverage is the $\lambda$-discounted
average coverage of its out-neighbors. Setting $\lambda = 1$ recovers the
undiscounted score, in which a token found $j$ hops away and a token found
at the root contribute equally; for $\lambda < 1$, a token found $j$ hops
away instead contributes $\lambda^j$, so the score also rewards
\emph{proximity} to a token, not only reachability. The overall objective
is the mean coverage of the mapped concepts at the full horizon $D$:
\[
F_D(T) = \frac{1}{|M|} \sum_{c \in M} S_D(c, T), \qquad F_D(\varnothing) = 0,\ \ 0 \le F_D(T) \le 1.
\]

This objective has a property that makes it tractable at scale: $F_D$ is
normalized, monotone, and submodular in $T$ (the marginal gain of adding a
token never increases as more tokens are already selected), which follows
from writing $F_D$ as a positive average, over random depth-$D$ walks from
each $c \in M$, of the indicator that the walk hits $T$. Submodularity buys
two things. First, the classical greedy approximation guarantee,
\begin{equation}
F_D(T^{\text{greedy}}_k) \ge \Big(1 - \tfrac{1}{e}\Big) \max_{|T| \le k} F_D(T).
\label{eq:guarantee}
\end{equation}
Second, and more important for a candidate pool of tens of thousands of
nodes, it licenses \emph{lazy} greedy: each candidate's previously computed
marginal gain is a valid upper bound on its current gain, so at every step
it suffices to re-verify only the best stale candidate against the
next-best stale bound, instead of recomputing every remaining candidate.
Algorithm~\ref{alg:select-tokens} states this loop precisely; it calls
\textsc{CandidateDelta} (Algorithm~\ref{alg:candidate-delta}), which
computes one candidate's exact marginal gain by propagating a sparse
\emph{delta} backward through the incoming-edge structure -- adding a
candidate $t$ only ever changes $S_h$ at $t$ itself and at vertices that can
reach $t$ within $D$ hops in reverse -- instead of recomputing the full
$O(D(|V|+|E|))$ recursion of Eq.~\eqref{eq:coverage} from scratch for every
candidate at every step.

\begin{algorithm}[h]
\caption{\textsc{SelectTokens} -- lazy-greedy token selection}
\label{alg:select-tokens}
\begin{algorithmic}[1]
\Require candidate pool $V_D$, mapped concepts $M$, horizon $D$, budget $k$, penalty $\lambda$
\Ensure token set $T$ with $|T| = k$ (or fewer if candidates run out)
\State $T \gets \varnothing$;\ $S \gets \mathbf{0}$ \Comment{$S_h(u,T){=}0$ for all $u,h$ when $T=\varnothing$}
\ForAll{$t \in V_D$} \Comment{initial seeding pass -- exact gains, parallelizable}
    \State $(\text{gain}, \_) \gets \Call{CandidateDelta}{t}$
    \State push $(t, \text{gain})$ onto max-heap $H$
\EndFor
\While{$|T| < k$ and $H \neq \varnothing$}
    \Repeat
        \State $t \gets$ pop top of $H$
        \State $(\text{gain}, \delta) \gets \Call{CandidateDelta}{t}$ \Comment{re-verify exactly, not from a stale bound}
        \State $b \gets$ gain bound now at the top of $H$ (or $-\infty$ if empty)
    \Until{$\text{gain} \ge b$} \Comment{submodularity: a stale bound never overstates a candidate's true gain}
    \State $T \gets T \cup \{t\}$;\ \Call{Commit}{$t, \delta$}
\EndWhile
\State \Return $T$
\end{algorithmic}
\end{algorithm}

\begin{algorithm}[h]
\caption{\textsc{CandidateDelta} and \textsc{Commit} -- sparse incremental update}
\label{alg:candidate-delta}
\begin{algorithmic}[1]
\Function{CandidateDelta}{$t$}
    \State $\delta_0 \gets \{\, t \mapsto 1 - S_0(t) \,\}$ \Comment{only $t$ changes at horizon $0$}
    \For{$h = 1, \dots, D$}
        \State $\delta_h \gets \{\}$
        \ForAll{$(v, d_v) \in \delta_{h-1}$}
            \ForAll{$u \in \mathrm{In}(v)$ with $u \notin T \cup \{t\}$}
                \State $\delta_h[u] \gets \delta_h[u] + \lambda\, d_v / d^+(u)$ \Comment{Eq.~\eqref{eq:coverage}, propagated backward}
            \EndFor
        \EndFor
        \State $\delta_h[t] \gets 1 - S_h(t)$ \Comment{$t$ is absorbing at every horizon}
    \EndFor
    \State $\text{gain} \gets \frac{1}{|M|}\sum_{c \in M} \delta_D[c]$
    \State \Return $(\text{gain}, \delta)$
\EndFunction
\Function{Commit}{$t$, $\delta$}
    \For{$h = 0, \dots, D$}
        \ForAll{$(u, d_u) \in \delta_h$}
            \State $S_h(u) \gets S_h(u) + d_u$
        \EndFor
    \EndFor
    \State $T \gets T \cup \{t\}$
\EndFunction
\end{algorithmic}
\end{algorithm}

\subsubsection{Precomputation and reuse across a sweep}
\label{sec:method:precompute}

Two pieces of Algorithms~\ref{alg:select-tokens}--\ref{alg:candidate-delta}
depend only on the graph $G$, never on $T$ or on which candidate is being
scored: the row-normalized transition matrix $A$ that Eq.~\eqref{eq:coverage}
multiplies against (with $\lambda$ already folded in), and the
per-node out-edge lists used by the tokenizer
(Section~\ref{sec:method:tokenizer}), pre-sorted so that non-\texttt{IS\_A}
relationships are visited before \texttt{IS\_A} ones. Because a full
evaluation sweep re-runs selection for every $(\lambda, k)$ pair on the
\emph{same} graph, both structures are built exactly once per graph and
shared for the whole sweep -- via a process-pool initializer when the
initial seeding pass of Algorithm~\ref{alg:select-tokens} is parallelized --
rather than being rebuilt inside every task. This removes a redundant
$O(|V|+|E|)$ rebuild from what would otherwise be millions of repeated
calls. Only the seeding pass is parallelized this way: once the lazy loop
itself is running, an accepted token typically needs only $1$--$7$
re-evaluations of the heap's top candidate before it is certified as the
true maximizer, so the loop is kept serial -- at that scale, the overhead
of shipping work to a worker process would exceed the work saved.

\subsection{Tokenizer}
\label{sec:method:tokenizer}

\begin{quote}\small\itshape
\begin{itemize}
\item tokenizer (say why flatten is not good, add illustration, resulting context trees)
\end{itemize}
\end{quote}

Selecting $T$ answers "which concepts become tokens." A separate question
is how a mapped concept $c \in M$ should be \emph{represented} once $T$ is
fixed. The simplest option -- the flat, unordered set of tokens reachable
from $c$ -- throws away exactly the information that non-\texttt{IS\_A}
relationships were designed to carry: if a concept has two independent
\emph{finding site} occurrences, each qualified by its own
\emph{laterality}, a flat token set $\{\text{kidney}, \text{adrenal gland},
\text{left}, \text{right}\}$ can no longer tell which laterality belongs to
which site. \doubt{keeping this exact kidney/adrenal example from your PDF
-- confirm, or swap for a real HUG-mapped concept once we pick an
illustrative one for the figure below.}

We avoid this by recursively expanding each mapped concept $c$ into a
\emph{context tree} instead. \texttt{IS\_A} relationships are transparent
and continue the search in the current context; every other relationship
type opens a fresh subcontext, so that two independent occurrences of the
same non-\texttt{IS\_A} relationship type are kept apart. Algorithm~\ref{alg:expand}
states this precisely; two implementation details are worth noting
alongside it. First, a node's outgoing edges are visited non-\texttt{IS\_A}
first, so that when the same subcontext key $(r, v)$ is reached twice at
different depths, the shallower (more direct) occurrence is what gets kept
(line~9's "keep more direct" rule in the underlying implementation).
Second, a token located exactly at depth $D$ is still accepted, since the
token check (line~2) happens before the depth-budget check (line~3).

\begin{algorithm}[h]
\caption{\textsc{Expand} -- context-tree construction, called as \textsc{Expand}$(c, q_0, 0)$ for each $c \in M$}
\label{alg:expand}
\begin{algorithmic}[1]
\Function{Expand}{$u$, $q$, $d$}
    \If{$u \in T$}
        \State attach $u$ as a token of context $q$ at depth $d$; \Return
    \ElsIf{$d = D$ or $u$ has no outgoing edges}
        \State mark $q$ \emph{uncovered} on this branch; \Return
    \Else
        \ForAll{outgoing edges $(u, r, v)$ of $u$}
            \If{$r = \texttt{IS\_A}$}
                \State \Call{Expand}{$v$, $q$, $d+1$} \Comment{transparent: stay in the same context}
            \Else
                \State $q' \gets$ subcontext of $q$ keyed by $(r, v)$, creating it if new
                \If{$(r,v)$ was new for $q$}
                    \State \Call{Expand}{$v$, $q'$, $d+1$}
                \EndIf
            \EndIf
        \EndFor
    \EndIf
\EndFunction
\end{algorithmic}
\end{algorithm}

The resulting tree keeps only $T$ as its vocabulary; its subcontext
structure records exactly which relationship chain each token was reached
through, which is what preserves the kidney/left vs.\ adrenal gland/right
distinction above.

\begin{figure}[h]
\centering
\fbox{\parbox{0.85\textwidth}{\centering
[FIGURE PLACEHOLDER -- side-by-side: flat token set vs. context tree for a
real mapped concept with $\ge 2$ non-\texttt{IS\_A} relationship occurrences,
e.g. the kidney/adrenal-gland example, or a concept pulled directly from the
HUG-mapped data.]}}
\caption{Context-tree tokenization example (to be replaced with an actual
figure).}
\label{fig:context-tree}
\end{figure}

\subsection{Evaluation metrics}
\label{sec:method:metrics}

\begin{quote}\small\itshape
\begin{itemize}
\item metrics used and how they are computed,
\item 3.1: Efficientcy: conciseness, tree\_complexity
\item 3.2: distance\_score: Fidelity
\item 3.3: discrimitive power: uniqueness\_entropy, unique\_rate
\item 3.4: semantic content: semantic\_coverage, uncovered\_rate, unk\_branche\_rate
\item and also say either they are hogher or lower better (interpretation)
\end{itemize}
\end{quote}

With a fixed $T$ and context trees for all of $M$ in hand, we can now ask
how good this representation actually is. All eight metrics below are
computed from the \emph{realized} context trees, as opposed to $F_D(T)$,
which is only the objective used to \emph{select} $T$ in the first place
(Section~\ref{sec:method:selection}). Table~\ref{tab:metrics} groups them
into four families and states, for each, which direction is better.

\begin{table}[h]
\centering
\small
\begin{tabular}{p{2.3cm}p{2.0cm}p{1.0cm}p{5.5cm}}
\toprule
\textbf{Metric} & \textbf{Family} & \textbf{Better} & \textbf{Definition} \\
\midrule
Conciseness & Efficiency & lower & Mean number of token leaves per concept's context tree. \\
Tree complexity & Efficiency & lower & Mean number of contexts (root plus subcontexts) per tree. \\
Distance score & Fidelity & higher & $1 - \overline{d}/(D{+}1)$, where $\overline{d}$ is the mean hop-distance at which tokens are found (uncovered branches counted at $D{+}1$). $1.0$ means every branch found a token immediately. \\
Uniqueness entropy & Discriminative power & higher & Normalized Shannon entropy of context-tree signatures across $M$; low when many concepts collapse onto the same tree shape. \\
Unique rate & Discriminative power & higher & Per-concept decomposition of the above: fraction of concepts whose signature is not shared by any other concept. \\
Exact rate & Discriminative power & higher & Fraction of mapped concepts that are themselves selected as a token (a $0$-hop, exact representation); maximally discriminable by construction. \\
Semantic coverage & Semantic content & higher & $F_D(T)$ itself, always reported at $\lambda=1$ regardless of the $\lambda$ used at \emph{selection} time (see \doubt{confirm this evaluation convention}). \\
Uncovered rate & Semantic content & lower & Fraction of concepts with \emph{zero} tokens anywhere in their tree. \\
Unk rate & Semantic content & lower & Fraction of concepts with \emph{at least one} uncovered branch; by construction, uncovered rate $\le$ unk rate. \\
\bottomrule
\end{tabular}
\caption{The eight evaluation metrics, grouped by family, with direction of improvement and definition.}
\label{tab:metrics}
\end{table}

\subsection{Baselines}
\label{sec:method:baselines}

\begin{quote}\small\itshape
\begin{itemize}
\item Chosen baseline for token selection (all DAG compatible)
\item 4.0: k random
\item 4.1: degree based: Most children, Highest degree, Highest degree with D <= 1
\item 4.2: centrality based: Closeness centrality, Page rank, Personalized page rank
\item 4.3 ours, and also our ablated version: Discrete set cover
\end{itemize}
\end{quote}

We compare our method against seven baselines from three families, plus one
ablation of our own method, summarized in Table~\ref{tab:baselines}. All of
them are chosen to be well defined on a DAG; \notetoself{per your feedback,
eigenvector centrality has been removed from this list entirely -- its
Perron--Frobenius uniqueness guarantee requires strong connectivity, which a
DAG cannot have, so it is no longer part of this draft even as a
caveated option. It remains a real baseline in the codebase/notebooks; this
is a paper-scope decision, not a code change.} Every baseline ranks the full
candidate pool $V_D$ once and is then truncated to the top-$k$ prefix of
that ranking, exactly as our own method's selection history is truncated --
so every method in Section~\ref{sec:results} is compared at matched $k$.

\begin{table}[h]
\centering
\small
\begin{tabular}{p{2.2cm}p{1.7cm}p{4.9cm}p{2.0cm}}
\toprule
\textbf{Baseline} & \textbf{Family} & \textbf{Rule} & \textbf{DAG note} \\
\midrule
$k$-random & Random & Independent uniform sample of $k$ candidates without replacement, repeated over seeds and averaged (no single ranking to truncate). & Trivial \\
Highest degree & Degree-based & Number of distinct predecessors within $D$ hops, over every relationship type. & Compatible \\
Highest degree, $\Delta\le 1$ & Degree-based & Same score, restricted to direct ($1$-hop) predecessors. & Compatible \\
Most children & Degree-based & Same $D$-hop predecessor count, restricted to the \texttt{IS\_A}-only subgraph. & Compatible \\
PageRank & Centrality & Stationary distribution of a random walk with restart probability $1-\alpha$, $\alpha=0.85$ \citeneeded{Brin \& Page, PageRank}. & Compatible \\
Personalized PageRank & Centrality & Same recursion, restarting uniformly on $M$ instead of on all of $V$ -- the one $M$-aware centrality baseline \citeneeded{personalized / topic-sensitive PageRank, e.g. Haveliwala}. & Compatible \\
Closeness centrality & Centrality & Reciprocal of the mean shortest-path distance \emph{to} a candidate from every node that can reach it (in-degree-oriented, Wasserman--Faust improved) \citeneeded{Freeman, closeness centrality}. & Compatible \\
Discrete set-cover (ablation) & Ours, ablated & Same lazy/CELF greedy mechanism and $(1{-}1/e)$ guarantee (Eq.~\eqref{eq:guarantee}) as our method, but $S_h(u,T)$ replaced by a binary "within $D$ hops of some token" indicator -- classical maximum-coverage greedy. & Compatible \\
\bottomrule
\end{tabular}
\caption{Baselines compared against our method, grouped by family. Eigenvector centrality is excluded (see note below).}
\label{tab:baselines}
\end{table}

The discrete set-cover ablation is the one entry that is not a baseline in
the usual sense: it isolates what the soft, $\lambda$-decayed, multi-hop
coverage of Eq.~\eqref{eq:coverage} buys over textbook hard-coverage greedy,
using the \emph{same} selection mechanism, so any difference in
Section~\ref{sec:results} can be attributed to the coverage formulation
itself rather than to the greedy procedure around it.

\subsection{Expert acceptability evaluation}
\label{sec:method:expert-eval}

\begin{quote}\small\itshape
\begin{itemize}
\item evaluation of accaptability of tokenized version versus ooriginal concept, by clinical expert.
\end{itemize}
\end{quote}

The metrics above are all automatic and intrinsic to the graph; they do not
ask whether a person would actually recognize the tokenized representation
as faithful to the original concept. We close this gap with a clinical
expert evaluation: for a sample of mapped concepts, the expert is shown the
original concept alongside its context-tree tokenization
(Section~\ref{sec:method:tokenizer}) and rates whether the tokenized version
still conveys an acceptable representation of the original concept's
meaning. \notetoself{per your feedback, this is now firm in-paper content
(no longer "(maybe)"), and it is intrinsic quality evaluation, distinct from
the downstream/assistive-clinical-decision-task validation named in
Section~\ref{sec:limitations} -- the two are not the same thing and this
draft keeps them separate.} \doubt{still need from you: how many
concepts/experts, what rating scale (binary acceptable/not, Likert,
free-text justification?), and whether inter-rater agreement is reported if
there is more than one expert -- these design choices affect how this
subsection should be written up and what Section~\ref{sec:results} needs to
report alongside the automatic metrics.}
